\documentclass[11pt,a4paper]{article}
\usepackage[utf8]{inputenc}
\usepackage[margin=2cm]{geometry}
\usepackage{amsmath,amssymb}
\usepackage{enumitem}
\usepackage{fancyhdr}
\usepackage{tabularx}
\usepackage{booktabs}

\setlength{\headheight}{14pt}
\pagestyle{fancy}
\fancyhf{}
\lhead{\textbf{MA1101 Matematika I --- Kuis Singkat (20 Menit)}}
\rhead{\textbf{Semester 1 TA 2026/2027}}
\rfoot{\thepage}

\begin{document}

% =========================================================================
% PAKET A
% =========================================================================
\begin{center}
    {\large \textbf{FAKULTAS MATEMATIKA DAN ILMU PENGETAHUAN ALAM}}\\
    {\large \textbf{INSTITUT TEKNOLOGI BANDUNG}}\\[4pt]
    {\large \textbf{KUIS KELAS MA1101 MATEMATIKA I}}\\[2pt]
    \textbf{Materi: Subbab 2.1--2.7, 2.9, 3.1--3.3 \quad | \quad Durasi: 20 Menit \quad | \quad Sifat: Tutup Buku \& Tanpa Kalkulator}
\end{center}

\vspace{2pt}
\noindent\fbox{%
\begin{minipage}{0.98\textwidth}
\begin{tabular}{ll@{\hspace{2cm}}ll}
    \textbf{Nama Mahasiswa} & : \dotfill & \textbf{Kelas} & : K10 / K35 / K40 \\
    \textbf{NIM}            & : \dotfill & \textbf{Paket} & : \textbf{A} \\
\end{tabular}
\end{minipage}%
}

\vspace{10pt}
\noindent\textbf{Petunjuk:} Jawablah setiap pertanyaan di bawah ini dengan menyertakan langkah pengerjaan yang jelas dan terstruktur. Tuliskan jawaban Anda pada lembar yang disediakan.

\vspace{8pt}
\hrule
\vspace{10pt}

% SOAL 1 (PAKET A)
\noindent\textbf{Soal 1 [Bobot: 45 Poin] (Definisi Turunan, Garis Singgung, Turunan Implisit, \& Hampiran)}
\begin{enumerate}[label=(\alph*), leftmargin=24pt]
    \item \textbf{[20 Poin]} Diberikan fungsi $f(x) = x^2 + 2x$. Tentukan $f'(2)$ dan tentukan hampiran linear dari $f$ di sekitar $x=2$.
    \item \textbf{[25 Poin]} Diberikan kurva implisit yang didefinisikan oleh persamaan:
    \[
    x^2 y = 4
    \]
    \begin{enumerate}[label=(\roman*)]
        \item Gunakan turunan implisit untuk menentukan nilai $\dfrac{dy}{dx}$ di titik $(2, 1)$.
        \item Tentukan persamaan garis singgung pada kurva tersebut di titik $(2, 1)$.
    \end{enumerate}
\end{enumerate}

\vspace{14pt}

% SOAL 2 (PAKET A)
\noindent\textbf{Soal 2 [Bobot: 55 Poin] (Ekstrem, Kemonotonan, Kecekungan, \& Titik Belok)}
\vspace{4pt}

\noindent Diberikan fungsi $f(x) = x^3 - 3x$ yang terdefinisi pada selang tutup $[-2, 2]$.
\begin{enumerate}[label=(\alph*), leftmargin=24pt]
    \item \textbf{[20 Poin]} Tentukan selang di mana $f$ monoton naik dan selang di mana $f$ monoton turun. Tentukan pula semua titik ekstrem lokal beserta jenisnya (maksimum lokal atau minimum lokal).
    \item \textbf{[15 Poin]} Tentukan selang di mana grafik $f$ cekung ke atas dan selang di mana grafik $f$ cekung ke bawah, serta tentukan koordinat titik beloknya.
    \item \textbf{[20 Poin]} Tentukan semua titik kritis dari $f$ pada selang $[-2, 2]$, kemudian tentukan nilai \textbf{maksimum global} dan nilai \textbf{minimum global} dari $f$ pada selang tersebut.
\end{enumerate}

\newpage

% =========================================================================
% PAKET B
% =========================================================================
\begin{center}
    {\large \textbf{FAKULTAS MATEMATIKA DAN ILMU PENGETAHUAN ALAM}}\\
    {\large \textbf{INSTITUT TEKNOLOGI BANDUNG}}\\[4pt]
    {\large \textbf{KUIS KELAS MA1101 MATEMATIKA I}}\\[2pt]
    \textbf{Materi: Subbab 2.1--2.7, 2.9, 3.1--3.3 \quad | \quad Durasi: 20 Menit \quad | \quad Sifat: Tutup Buku \& Tanpa Kalkulator}
\end{center}

\vspace{2pt}
\noindent\fbox{%
\begin{minipage}{0.98\textwidth}
\begin{tabular}{ll@{\hspace{2cm}}ll}
    \textbf{Nama Mahasiswa} & : \dotfill & \textbf{Kelas} & : K10 / K35 / K40 \\
    \textbf{NIM}            & : \dotfill & \textbf{Paket} & : \textbf{B} \\
\end{tabular}
\end{minipage}%
}

\vspace{10pt}
\noindent\textbf{Petunjuk:} Jawablah setiap pertanyaan di bawah ini dengan menyertakan langkah pengerjaan yang jelas dan terstruktur. Tuliskan jawaban Anda pada lembar yang disediakan.

\vspace{8pt}
\hrule
\vspace{10pt}

% SOAL 1 (PAKET B)
\noindent\textbf{Soal 1 [Bobot: 45 Poin] (Definisi Turunan, Garis Singgung, Turunan Implisit, \& Hampiran)}
\begin{enumerate}[label=(\alph*), leftmargin=24pt]
    \item \textbf{[20 Poin]} Diberikan fungsi $g(x) = x^2 + 3x$. Tentukan $g'(1)$ dan tentukan hampiran linear dari $g$ di sekitar $x=1$.
    
    \item \textbf{[25 Poin]} Diberikan kurva implisit yang didefinisikan oleh persamaan:
    \[
    x^2y = 9
    \]
    \begin{enumerate}[label=(\roman*)]
        \item Gunakan turunan implisit untuk menentukan nilai $\dfrac{dy}{dx}$ di titik $(3, 1)$.
        \item Tentukan persamaan garis singgung pada kurva tersebut di titik $(3, 1)$.

    \end{enumerate}
\end{enumerate}

\vspace{14pt}

% SOAL 2 (PAKET B)
\noindent\textbf{Soal 2 [Bobot: 55 Poin] (Ekstrem, Kemonotonan, Kecekungan, \& Titik Belok)}
\vspace{4pt}

\noindent Diberikan fungsi $g(x) = x^3 - 12x$ yang terdefinisi pada selang tutup $[-3, 3]$.
\begin{enumerate}[label=(\alph*), leftmargin=24pt]
    \item \textbf{[20 Poin]} Tentukan selang di mana $g$ monoton naik dan selang di mana $g$ monoton turun. Tentukan pula semua titik ekstrem lokal beserta jenisnya (maksimum lokal atau minimum lokal).
    \item \textbf{[15 Poin]} Tentukan selang di mana grafik $g$ cekung ke atas dan selang di mana grafik $g$ cekung ke bawah, serta tentukan koordinat titik beloknya.
    \item \textbf{[20 Poin]} Tentukan semua titik kritis dari $g$ pada selang $[-3, 3]$, kemudian tentukan nilai \textbf{maksimum global} dan nilai \textbf{minimum global} dari $g$ pada selang tersebut.
\end{enumerate}

\end{document}
